\chapter*{Abstract}
\markboth{Abstract}{Abstract}

Seagrass meadows bury organic carbon at rates that rival tropical forests,
yet the permanence of that store under a warming ocean is poorly
constrained. This thesis asks how eelgrass (\emph{Zostera marina})
meadows in a temperate estuary respond to marine heatwaves, and whether
the carbon they have accumulated over centuries remains in the sediment
once the canopy thins.

Combining six years of monthly shoot censuses across twelve meadows,
in situ warming chambers and \(^{210}\)Pb dated sediment cores,
I show that shoot density declines non-linearly once the cumulative
heat dose exceeds roughly \(14\) degree-heating-days, a threshold crossed
in three of the six summers studied. Below the threshold meadows recover
within a single growing season; above it, recovery takes two to four
years and the upper ten centimetres of sediment lose between \(18\) and
\(31\) per cent of their organic carbon to remineralisation.

These results suggest that blue-carbon accounting which assumes a
permanent sediment store overestimates the climate benefit of temperate
meadows by up to a quarter under current heatwave frequencies, and they
identify the summer heat dose as a practical early-warning indicator for
managers.

\cleardoublepage
